% Figure 1.1 --- L'arc de l'ancrage : trois générations d'approches
\begin{tikzpicture}

\node[fbleu, anchor=north west, align=left, text width=5.1cm, inner sep=7pt]
      at (0,0) (a)
     {{\bfseries\color{figbleu} 1. Classification (axe 1)}\\[3pt]
      \emph{Mécanisme} : un modèle affiné range un texte dans une catégorie
      ATT\&CK (TRAM) ou en extrait un graphe (AttacKG).\\[3pt]
      \emph{Chiffre} : GPT-4 $\approx$ 0,64 F1, Llama-3-8B $\approx$ 0,16
      (CTIBench).\\[3pt]
      \emph{Limite} : couverture partielle, plafond structurel.};

\node[fbord, anchor=north west, align=left, text width=5.1cm, inner sep=7pt]
      at (6.4cm,0.55cm) (b)
     {{\bfseries\color{figbordeaux} 2. Récupération --- RAG (axe 2)}\\[3pt]
      \emph{Mécanisme} : le modèle raisonne sur des passages récupérés par
      similarité vectorielle.\\[3pt]
      \emph{Chiffre} : 5 textes piégés dans 2,6 M de documents
      $\Rightarrow$ 90 \% d'attaques réussies (PoisonedRAG).\\[3pt]
      \emph{Limite} : voisin $\neq$ pertinent ; base empoisonnable.};

\node[fvert, anchor=north west, align=left, text width=5.1cm, inner sep=7pt]
      at (12.8cm,0) (c)
     {{\bfseries\color{figvert} 3. Outils --- MCP (axe 3)}\\[3pt]
      \emph{Mécanisme} : le modèle exécute une requête structurée sur une base
      formelle indexée.\\[3pt]
      \emph{Chiffre} : un modèle outillé de 6,7 G bat GPT-3 (175 G) $\times$3
      en arithmétique (Toolformer).\\[3pt]
      \emph{Limite} : le risque passe du texte à l'action (axe 4).};

\draw[fgrosse] ($(a.east)+(0,-0.3cm)$) -- ($(b.west)+(0,-0.85cm)$);
\draw[fgrosse] ($(b.east)+(0,-0.85cm)$) -- ($(c.west)+(0,-0.3cm)$);

\node[below=3pt of a.south, font=\scriptsize, text=black!55] {$\sim$2021--2024};
\node[below=3pt of b.south, font=\scriptsize, text=black!55] {2024--2025};
\node[below=3pt of c.south, font=\scriptsize, text=black!55] {2025--2026};

\end{tikzpicture}
